\chapter{Ipotesi e criteri di falsificazione}\label{chap:hypotheses}


\noindent Questo capitolo documenta il processo che ha condotto dalle osservazioni iniziali alla
formulazione delle ipotesi. La sua inclusione risponde a una ragione metodologica: il lavoro non è
consistito nell'applicare una tecnica nota a un sistema noto, ma nel ricostruire il comportamento
di un sistema di cui non erano disponibili i sorgenti. Le ipotesi hanno quindi svolto una funzione
euristica, e la loro falsificazione è parte del risultato. Il capitolo ricostruisce il processo di
analisi, enuncia le ipotesi primarie con la loro derivazione e il criterio di falsificazione ---
senza anticiparne l'esito, che è discusso insieme ai risultati --- e raccoglie separatamente le
domande esplorative. Il protocollo di verifica e la metodologia statistica, che delle ipotesi
costituiscono il banco di prova, sono esposti nel capitolo sulla valutazione sperimentale.

% =======================================================================================
\section{Come sono nate le ipotesi}\label{sec:analysis-process}

L'analisi si è svolta in quattro passi, ciascuno dei quali ha determinato il successivo.

\begin{enumerate}[leftmargin=1.6em]
    \item \textbf{Disassemblaggio.} Il sistema è distribuito senza sorgenti. La ricostruzione
    della logica è avvenuta per disassemblaggio del bytecode, con particolare attenzione alle
    procedure che consumano l'insieme delle regole.
    \item \textbf{Individuazione del doppio ruolo.} È emerso che una RFD partecipa a due
    procedure con requisiti opposti (\autoref{sec:dual-role}). Questa osservazione ha trasformato
    il problema da ``selezionare le regole utili'' a ``selezionare le regole utili sapendo che
    \emph{utile} ha due significati incompatibili''.
    \item \textbf{Derivazione delle condizioni.} Dal comportamento delle due procedure sono state
    ricavate le condizioni \Vone{}, \Vtwo{}, \Vthree{} e il teorema negativo che ne segue
    (\autoref{chap:theory}).
    \item \textbf{Verifica sperimentale.} Le condizioni sufficienti sono state verificate su
    stati reali; le strategie euristiche sono state misurate su più famiglie di dataset, con
    particolare attenzione a \emph{dove} esse funzionano e dove no.
\end{enumerate}

Un elemento del processo merita di essere esplicitato, perché ha modificato il corso del lavoro.
L'analisi del passo 3 ha prodotto un risultato negativo, il teorema della
\autoref{sec:negative-theorem}, che esclude la possibilità di raggiungere l'obiettivo di
riduzione mantenendo la garanzia. Una parte non trascurabile del lavoro è consistita nel
\emph{verificare} che quel risultato negativo non fosse un artefatto
dell'implementazione, prima di accettarlo come limite del problema
(\autoref{sec:maxextra}).

\paragraph{Origine delle ipotesi.} Il lavoro è esplorativo: le ipotesi non precedono l'analisi ma
nascono da essa, e la circostanza va dichiarata perché incide sul valore probatorio dei test. La
\autoref{tab:hypotheses} distingue perciò, per ogni ipotesi, l'\emph{origine}: le ipotesi
\emph{derivate} discendono da una proprietà stabilita prima di osservare i dati sperimentali ---
una condizione dimostrata, una misura di struttura, l'analisi del costo --- e la loro previsione è
argomentabile a priori; le ipotesi \emph{esplorative} sono state formulate durante l'analisi, e il
loro esito va inteso come indicazione per ulteriori verifiche e non come conferma indipendente.
Nessuna delle ipotesi primarie è stata formulata dopo aver visto l'esito dell'esperimento che la
verifica; la distinzione è mantenuta anche nel capitolo dei risultati, dove gli esiti sono
riportati per ciascuna.

% =======================================================================================
\section{Ipotesi primarie}\label{sec:hypotheses-list}

Le ipotesi primarie sono sette e corrispondono alle domande che il lavoro pone: che cosa si può
potare con garanzia, che cosa si può preservare senza garanzia, se esista una leva di qualità, se
tale leva generalizzi, da che cosa dipenda il guadagno quando si manifesta, e che rapporto vi sia
fra riduzione e tempo. La \autoref{tab:hypotheses} le enuncia con la previsione e il criterio di
falsificazione; l'esito non è qui riportato.

\begin{table}[htbp]
\centering
\caption{Ipotesi primarie, previsione, criterio di falsificazione e sezione in cui ciascuna è
risolta. L'origine --- derivata da una proprietà stabilita prima della misura, oppure formulata
durante l'analisi --- è dichiarata nella derivazione che segue alla tabella.}
\label{tab:hypotheses}
\footnotesize
\begin{tabular}{@{}cp{3.3cm}p{3.3cm}p{2.3cm}c@{}}
\toprule
\# & ipotesi & previsione & falsificazione & risolta in \\
\midrule
H1 & \str{score-argmin} con veti protetti & riduzione $\ge \SI{30}{\percent}$ con output identico & riduzione inferiore alla soglia & \ref{sec:negative-results} \\
H6 & \str{veto-cover} su tutte le famiglie & output identico alla configurazione di partenza & una sola configurazione con output diverso & \ref{sec:res-vetocover} \\
H7 & verifica empirica di \Vtwo{} & zero discrepanze fra copertura prevista e verificata & una sola discrepanza & \ref{sec:res-vetocover} \\
H8 & \str{precision-priority} da solo & guadagno di qualità & nessun guadagno o perdita & \ref{sec:ablation} \\
H11 & \str{precision-priority} su Glass e Restaurant & guadagno anche fuori da Bridges & perdita su entrambe le famiglie & \ref{sec:res-boundary} \\
H14 & attribuzione del guadagno alla precisione dei cluster & un riordino casuale non produce lo stesso guadagno & il riordino casuale produce lo stesso guadagno & \ref{sec:negative-results} \\
H16 & proporzionalità fra riduzione e tempo & riduzioni maggiori producono tempi minori & una configurazione con riduzione maggiore e tempo non inferiore & \ref{sec:res-cost} \\
\bottomrule
\end{tabular}
\end{table}

\paragraph{H1 --- La strategia esatta lato generazione.} \str{score-argmin} preserva la generazione
dei candidati per costruzione, perché calcola staticamente l'argomento del minimo che il sistema
sceglierebbe. Ci si attende pertanto che, associata a una politica che tuteli integralmente il
potere di veto, essa produca un output identico a quello della configurazione di partenza, e che
la riduzione residua sia ancora significativa: la politica di tutela ripristina le regole
portatrici di veto, ma non quelle la cui rimozione è provata priva di costo dalla condizione
\Vtwo{}, e la classe \Vtwo{} vale fino a poco più di un quinto dell'insieme
(\autoref{tab:conditions}). La previsione è una riduzione non inferiore al \SI{30}{\percent}.
L'ipotesi è derivata: entrambe le premesse --- l'esattezza lato generazione e la dimensione della
classe coperta --- sono stabilite prima della misura.

\paragraph{H6 --- La classe garantita.} Se la condizione \Vtwo{} è corretta, la sua applicazione
deve essere invisibile al sistema: le regole rimosse hanno un veto implicato da una regola
sopravvissuta, e la proposizione della \autoref{prop:v2} assicura che nessun veto venga perso. La
previsione è che l'output coincida con quello della configurazione di partenza su \emph{tutte} le
configurazioni della griglia, non soltanto su quelle di una famiglia. È l'ipotesi con il criterio
di falsificazione più severo, perché una sola configurazione difforme la falsifica.

\paragraph{H7 --- Verifica indipendente della copertura.} La correttezza della condizione \Vtwo{}
è stabilita per via logica, ma la sua \emph{implementazione} potrebbe essere difettosa: la ricerca
dei copritori è un algoritmo, e un algoritmo può sbagliare. La previsione è che, controllando a
posteriori su stati reali del dataset che ogni regola rimossa abbia effettivamente un copritore
sopravvissuto che esercita il veto sugli stessi testimoni, non emerga alcuna discrepanza. La
verifica è indipendente perché non riusa il codice della condizione ma ricostruisce i testimoni.

\paragraph{H8 --- La leva di qualità.} L'analisi del sistema mostra che l'esito di
un'imputazione dipende dall'ordine in cui i cluster vengono tentati
(\autoref{sec:two-branches}), e che la posizione di una regola nell'ordine di scansione non
richiede la rimozione di alcunché. Ne segue la previsione che un riordino che disponga per prime
le regole con la precisione empirica più alta migliori il risultato senza toccare il potere di
veto. L'ipotesi è derivata dalla struttura delle due procedure, non da un'osservazione sui
risultati.

\paragraph{H11 --- Generalizzazione della leva di qualità.} Se il meccanismo del guadagno è la
composizione dei candidati e non una proprietà di Bridges, lo stesso riordino deve produrre un
guadagno anche sulle altre famiglie. La previsione è che il guadagno si manifesti su Glass e su
Restaurant; la falsificazione è una perdita su entrambe. L'ipotesi è deliberatamente severa: è la
generalizzazione a decidere se il guadagno osservato su una famiglia sia un fenomeno del sistema o
un accidente del dataset.

\paragraph{H14 --- Attribuzione del guadagno.} Il riordino per priorità di precisione impiega un
segnale --- la precisione empirica dei cluster --- e la spiegazione naturale del guadagno è che
quel segnale sia la causa. Se la spiegazione è corretta, allora un riordino casuale non deve
produrre lo stesso guadagno: la previsione è che il guadagno del riordino informato superi in modo
sistematico quello di permutazioni con seme fissato. Si tratta di un'ipotesi di \emph{confutazione
di una spiegazione} più che di verifica di una strategia, ed è classificata come esplorativa
perché formulata dopo aver osservato il guadagno di cui cerca la causa.

\paragraph{H16 --- Rapporto fra riduzione e tempo.} L'analisi del costo
(\autoref{sec:cost}) mostra che il tempo è assorbito dalla verifica dei veti, che esamina
l'insieme delle regole per ogni candidato: se il costo per candidato è proporzionale al numero di
regole, la riduzione dell'insieme deve tradursi in una riduzione del tempo. La struttura del
sistema suggerisce però una previsione opposta: il veto non si limita a costare, ma \emph{scarta}
i candidati, e un insieme di regole più povero può lasciare in gara più alternative e quindi
moltiplicare le verifiche. L'ipotesi assume la prima previsione; la seconda è la ragione per cui
il criterio di falsificazione è enunciato in forma stretta --- una sola configurazione con
riduzione maggiore e tempo non inferiore falsifica l'ipotesi --- e non come semplice assenza di
correlazione.

% =======================================================================================
\section{Domande esplorative}\label{sec:exploratory}

Le restanti ipotesi non hanno il carattere di previsione argomentabile a priori: riguardano
varianti di parametri, strategie alternative e proprietà di portata locale, e sono state
formulate nel corso dell'analisi. Sono raccolte separatamente, senza esito, perché il loro ruolo è
delimitare lo spazio delle spiegazioni e indicare che cosa non funziona, non confermare una
previsione. La \autoref{tab:exploratory} le elenca con la domanda a cui rispondono.

\begin{table}[htbp]
\centering
\caption{Domande esplorative, con la sezione in cui ciascuna è risolta. Sono formulate durante
l'analisi; il loro esito è riportato nella \autoref{sec:hypotheses-outcomes} e non ha valore di
conferma indipendente.}
\label{tab:exploratory}
\footnotesize
\begin{tabular}{@{}cp{7.6cm}c@{}}
\toprule
\# & domanda & risolta in \\
\midrule
H2  & \str{score-argmin} preserva la copertura anche su Glass e Restaurant? & \ref{sec:negative-results} \\
H3  & esiste una strategia di copertura esatta delle celle, e a costo nullo? & \ref{sec:negative-results} \\
H4  & la conversione delle regole di tipo \texttt{B} in tipo \texttt{D} produce un guadagno gratuito? & \ref{sec:d1} \\
H5  & il riordino delle righe del dataset produce un guadagno? & \ref{sec:negative-results} \\
H9  & riordino e potatura sommano i loro effetti? & \ref{sec:ablation} \\
H10 & una regolarizzazione della stima delle precisioni produce un guadagno stabile? & \ref{sec:negative-results} \\
H12 & la versione regolarizzata generalizza meglio delle altre? & \ref{sec:negative-results} \\
H13 & il riordino entro il cluster è una leva di qualità? & \ref{sec:negative-results} \\
H15 & l'ordinamento delle rimozioni sotto vincolo di budget è migliorabile? & \ref{sec:negative-results} \\
\bottomrule
\end{tabular}
\end{table}

% =======================================================================================
